\section{Vectors as derivatives and their transformation fules}

Allright we saw that basis vectors can be re-interpreted as partial derivative of the position
vector along the coordinate axis.\\
We can sort of extend this logic to the vectors themselves. We saw that figuring out vector 
components in a particular basis is kind of easy to do for individual vectors, but now we're going do it 
for vector fields (i.e. a vector for every point in space), \textbf{but for now, we're going to deal with a slightly 
different type of vector field, and that's a vector field along a curve, so basically we're defining 
a vector at every point along a curve.}\\

The way we imagine vector fields along a curve is considering the set of tangent vectors to the curves.
This gives us a vector field defined at every point on the curve.\\

A curve can be tought as a function that takes some input parameter $\lambda$ and outputs the position vector 
$\mathbf{R}$ that traces our curve in the space.\\
If we want to have a vector tangent to this curve, we just need to consider taking the derivative of the position 
vector with respect to its parameter:

\begin{equation}
    \frac{\dd \begingroup\color{violet}\vec{R}\endgroup}{\dd \begingroup\color{orange}\lambda\endgroup}
\end{equation}

Now you may think that finding out the components of this vector in a specific basis can be 
difficult but it is really not, if we make use of the multivariable chain rule and expand it out in the
coodinates we're working in.\\
For example, in the cartesian coordinates $\color{cyan}c^i$ and using the Einstein summation convention notation:

\begin{equation}
    \frac{\dd \begingroup\color{violet}\vec{R}\endgroup}{\dd \begingroup\color{orange}\lambda\endgroup} = 
    \frac{\dd \begingroup\color{cyan}c^i\endgroup}{\dd \begingroup\color{orange}\lambda\endgroup}
    \frac{\partial \begingroup\color{violet}\vec{R}\endgroup}{\partial \begingroup\color{cyan}c^i\endgroup}
\end{equation}

At the end of the day, considering doing it in cartesian $\mathbb{R^2}$ for simplicity, these are equivalent (individual vector on the top, and vector field along a curve on the bottom):

\begin{align*}
    \begingroup\color{orange}\vec{v}\endgroup &= \begingroup\color{cyan}v^1\endgroup \begingroup\color{blue}\vec{e_1}\endgroup + \begingroup\color{cyan}v^2\endgroup \begingroup\color{blue}\vec{e_2}\endgroup\\
    \frac{\dd \begingroup\color{violet}\vec{R}\endgroup}{\dd \begingroup\color{orange}\lambda\endgroup} &= \frac{\dd \begingroup\color{cyan}x\endgroup}{\dd \begingroup\color{orange}\lambda\endgroup}
    \frac{\partial \begingroup\color{violet}\vec{R}\endgroup}{\partial \begingroup\color{cyan}x\endgroup} + 
    \frac{\dd \begingroup\color{cyan}y\endgroup}{\dd \begingroup\color{orange}\lambda\endgroup}
    \frac{\partial \begingroup\color{violet}\vec{R}\endgroup}{\partial \begingroup\color{cyan}y\endgroup}
\end{align*}

The same can be applied to whatever coordinate system we're working in, so to give you an example, let's say we
have the circular curve of radius 2, centered at the origin in $\mathbb{R^2}$, defined by:

\begin{equation}
    \begingroup\color{violet}\vec{R}\endgroup (\begingroup\color{orange}\lambda\endgroup) = \left(\begingroup\color{cyan}x\endgroup(\begingroup\color{orange}\lambda\endgroup), \begingroup\color{cyan}y\endgroup(\begingroup\color{orange}\lambda\endgroup)\right) =
    \left(2 \cos\lambda, 2 \sin\lambda\right)
\end{equation}

The vector field along the curve is made up of the tangent vectors along the curve, that 
can be expanded as defined:

\begin{align*}
    \frac{\dd \begingroup\color{violet}\vec{R}\endgroup}{\dd \begingroup\color{orange}\lambda\endgroup} &= \frac{\dd \begingroup\color{cyan}x\endgroup}{\dd \begingroup\color{orange}\lambda\endgroup}
    \frac{\partial \begingroup\color{violet}\vec{R}\endgroup}{\partial \begingroup\color{cyan}x\endgroup} + 
    \frac{\dd \begingroup\color{cyan}y\endgroup}{\dd \begingroup\color{orange}\lambda\endgroup}
    \frac{\partial \begingroup\color{violet}\vec{R}\endgroup}{\partial \begingroup\color{cyan}y\endgroup}\\
    &= -2 \sin\lambda \frac{\partial \begingroup\color{violet}\vec{R}\endgroup}{\partial \begingroup\color{cyan}x\endgroup} + 
    2 \cos\lambda\frac{\partial \begingroup\color{violet}\vec{R}\endgroup}{\partial \begingroup\color{cyan}y\endgroup}
\end{align*}

The components can also be written, for the cartesian basis, as a vector column:

\begin{equation}
    \begin{bmatrix}
        -2\sin\lambda\\
        2 \cos\lambda
    \end{bmatrix}_{\frac{\partial}{\partial \begingroup\color{cyan}c^i\endgroup}}
\end{equation}

We can do the same in polar coordinates $r, \theta$, where the same curve is defined as:

\begin{equation}
    \begingroup\color{violet}\vec{R}\endgroup (\begingroup\color{orange}\lambda\endgroup)= \left(\begingroup\color{magenta}r\endgroup(\begingroup\color{orange}\lambda\endgroup), \begingroup\color{magenta}\theta\endgroup(\begingroup\color{orange}\lambda\endgroup)\right) =
    \left(2 , \lambda\right)
\end{equation}

\begin{align*}
    \frac{\dd \begingroup\color{violet}\vec{R}\endgroup}{\dd \begingroup\color{orange}\lambda\endgroup} &= \frac{\dd \begingroup\color{magenta}r\endgroup}{\dd \begingroup\color{orange}\lambda\endgroup}
    \frac{\partial \begingroup\color{violet}\vec{R}\endgroup}{\partial \begingroup\color{magenta}r\endgroup} + 
    \frac{\dd \begingroup\color{magenta}\theta\endgroup}{\dd \begingroup\color{orange}\lambda\endgroup}
    \frac{\partial \begingroup\color{violet}\vec{R}\endgroup}{\partial \begingroup\color{magenta}\theta\endgroup}\\
    &= 0 \frac{\partial \begingroup\color{violet}\vec{R}\endgroup}{\partial \begingroup\color{magenta}r\endgroup} + 
    1 \frac{\partial \begingroup\color{violet}\vec{R}\endgroup}{\partial \begingroup\color{magenta}\theta\endgroup}
\end{align*}

So:

\begin{equation}
    \begin{bmatrix}
        0\\
        1
    \end{bmatrix}_{\frac{\partial}{\partial \begingroup\color{magenta}p^i\endgroup}}
\end{equation}

Now, as you can already imagine, like we saw that for individual vectors, the components are contravariant (i.e. they transform the opposite way of the basis vectors), 
we're going to simply show that this is valid indeed in this vector fields case.\\

We already saw in the previous section, the transformation rules of the basis vectors using the Jacobain and Inverse Jacobian:

\begin{subequations}\label{eq:forward-backward}
\begin{empheq}[box=\widefbox]{align}
    \frac{\partial \begingroup\color{violet}\vec{R}\endgroup}{\partial \begingroup\color{magenta}p^i\endgroup} &= \frac{\partial \begingroup\color{violet}\vec{R}\endgroup}{\partial \begingroup\color{cyan}c^j\endgroup}\frac{\partial \begingroup\color{cyan}c^j\endgroup}{\partial \begingroup\color{magenta}p^i\endgroup} & \text{Forward Transform}\\
    \frac{\partial \begingroup\color{violet}\vec{R}\endgroup}{\partial \begingroup\color{cyan}c^i\endgroup} &= \frac{\partial \begingroup\color{violet}\vec{R}\endgroup}{\partial \begingroup\color{magenta}p^j\endgroup}\frac{\partial \begingroup\color{magenta}p^j\endgroup}{\partial \begingroup\color{cyan}c^i\endgroup} & \text{Backward Transform}
\end{empheq}
\end{subequations}

Now we can just use these and replace the basis vectors with their new basis transformation rule: 

\begin{equation}
    \frac{\dd \begingroup\color{violet}\vec{R}\endgroup}{\dd \begingroup\color{orange}\lambda\endgroup} = 
    \frac{\dd \begingroup\color{cyan}c^i\endgroup}{\dd \begingroup\color{orange}\lambda\endgroup}
    \frac{\partial \begingroup\color{violet}\vec{R}\endgroup}{\partial \begingroup\color{cyan}c^i\endgroup} = 
    \frac{\dd \begingroup\color{cyan}c^i\endgroup}{\dd \begingroup\color{orange}\lambda\endgroup}
    \left(\frac{\partial \begingroup\color{violet}\vec{R}\endgroup}{\partial \begingroup\color{magenta}p^j\endgroup}\frac{\partial \begingroup\color{magenta}p^j\endgroup}{\partial \begingroup\color{cyan}c^i\endgroup}\right) = 
    \left(\frac{\dd \begingroup\color{cyan}c^i\endgroup}{\dd \begingroup\color{orange}\lambda\endgroup} \frac{\partial \begingroup\color{magenta}p^j\endgroup}{\partial \begingroup\color{cyan}c^i\endgroup}\right)
    \frac{\partial \begingroup\color{violet}\vec{R}\endgroup}{\partial \begingroup\color{magenta}p^j\endgroup}
\end{equation}

This is a linear combination of the new basis vector $\color{magenta}p^j$, and since we know that we can write:

\begin{equation}
    \frac{\dd \begingroup\color{violet}\vec{R}\endgroup}{\dd \begingroup\color{orange}\lambda\endgroup} = 
    \frac{\dd \begingroup\color{magenta}p^i\endgroup}{\dd \begingroup\color{orange}\lambda\endgroup}
    \frac{\partial \begingroup\color{violet}\vec{R}\endgroup}{\partial \begingroup\color{magenta}p^i\endgroup}
\end{equation}

This means that:

\begin{equation}
    \frac{\dd \begingroup\color{magenta}p^i\endgroup}{\dd \begingroup\color{orange}\lambda\endgroup} = 
    \frac{\dd \begingroup\color{cyan}c^i\endgroup}{\dd \begingroup\color{orange}\lambda\endgroup} \frac{\partial \begingroup\color{magenta}p^j\endgroup}{\partial \begingroup\color{cyan}c^i\endgroup}
\end{equation}

Which shows that the vector field components transform using the backward transformation (inverse Jacobian)\\

We can actually test this out in the previous example right?
So we have the components being, in both coordinate systems:

\begin{align*}
    \begin{bmatrix}
        -2\sin\lambda\\
        2 \cos\lambda
    \end{bmatrix}_{\frac{\partial }{\partial \begingroup\color{cyan}c^i\endgroup}} &&
    \begin{bmatrix}
        0\\
        1
    \end{bmatrix}_{\frac{\partial }{\partial \begingroup\color{magenta}p^i\endgroup}}
\end{align*}

And we also know the Inverse Jacobian in the case of polar-to-cartesian coordinates, being:

\begin{equation}
    B = J^{-1} =
    \begin{bmatrix}
        x/r & y/r\\
        -y/r^2 & x/r^2
    \end{bmatrix} = 
    \begin{bmatrix}
        \cos\lambda & \sin\lambda\\
        -1/2\sin\lambda & 1/2\cos\lambda
    \end{bmatrix}
\end{equation}

Let's check that:

\begin{align*}
    B \begin{bmatrix}
        -2\sin\lambda\\
        2 \cos\lambda
    \end{bmatrix}_{\frac{\partial }{\partial \begingroup\color{cyan}c^i\endgroup}} &= 
     \begin{bmatrix}
        \cos\lambda & \sin\lambda\\
        -1/2\sin\lambda & 1/2\cos\lambda
    \end{bmatrix} 
    \begin{bmatrix}
        -2\sin\lambda\\
        2 \cos\lambda
    \end{bmatrix}_{\frac{\partial }{\partial \begingroup\color{cyan}c^i\endgroup}}\\ 
    &= \begin{bmatrix}
        -\cos\lambda 2\sin\lambda + \sin\lambda 2\cos\lambda\\
        -1/2 \sin\lambda (-2 \sin\lambda) + 1/2 \cos\lambda 2 \cos\lambda
    \end{bmatrix}_{\frac{\partial}{\partial \begingroup\color{magenta}p^i\endgroup}}\\
    & = \begin{bmatrix}
        0\\
        1
    \end{bmatrix}_{\frac{\partial}{\partial \begingroup\color{magenta}p^i\endgroup}}
\end{align*}

Same can be checked on the opposite direction applying the forward/jacobian transformation.\\

One last thing to mention before terminating this section is in terms of notation: you notice that 
the position vector appears in a lot of them, so what we're going to do from here onward, is get rid of the 
position vector symbol.\\
This might seem like a weird choice of notation, but from now on, for example, the $\mathbb{R}^2$ cartesian basis vectors will be 
written as:

\begin{align*}
    \begingroup\color{blue}\vec{e_x}\endgroup = \frac{\partial}{\partial \begingroup\color{cyan}x\endgroup}\\
    \begingroup\color{blue}\vec{e_y}\endgroup = \frac{\partial}{\partial \begingroup\color{cyan}y\endgroup}
\end{align*}

\textbf{where basically we consider basis vectors these partial derivative operators.}\\
If you think about this, it's not very weird at the end, as derivative operators are linear, as they can be scaled and added like vectors do, 
so they have both a magnitude and a direction.
They also obey the forward and backward transformation rules that we would expect vectors to obey.\\

The main logical reason to get rid of the position vector symbol is that, the position vector definition requires an origin to be set, and 
later on when we're going to deal with curved surfaces and manifolds, it's not easy to define an origin point.\\
So, using a position vector without an origin does not make a lot of sense, but we're still going to be able to use
vectors and vector fields as long as we have these derivative operators to work with.\\



